% =====================================================================
%  Emisja polowa: sesja z magnesem vs. sesje bez magnesu
%  Fragment do wstawienia w rozdział 5 (\input{...}).
%  Wymaga: graphicx, subcaption, pgfplots (compat>=1.17), booktabs, siunitx
%  Ścieżki: figures/ (zdjęcia), dane/ (CSV z pomiarów i analizy obrazu)
% =====================================================================
\providecommand{\sesjaM}{z~magnesem}
\providecommand{\sesjaB}{bez magnesu}
\providecommand{\figdir}{figures}
\providecommand{\datadir}{dane}
\pgfplotsset{/pgf/number format/use comma, /pgf/number format/1000 sep={}}
\providecommand{\TODO}[1]{\textcolor{red}{[TODO: #1]}}

\section{Wpływ magnesu ogniskującego na emisję polową}
\label{sec:emisja-magnes}

\subsection{Przebieg pomiarów}

Porównano dwie sesje pomiarowe, w~których przy stałym napięciu targetu
zwiększano napięcie $U$ co \SI{100}{V} i~rejestrowano prąd emisji
$I_e$ oraz zdjęcie ekranu luminescencyjnego:
\begin{itemize}
  \item \textbf{Sesja B} (\sesjaB) -- napięcie targetu \SI{2000}{V},
        $U=\SIrange{500}{1500}{V}$, prąd w~zakresie
        \SIrange{2}{22}{\micro A};
        zdjęcia dla $U=\SIrange{700}{1100}{V}$ (arkusz \emph{Book1},
        seria~1). Przed tą sesją źródło zostało odgazowane przez
        wygrzanie w~temperaturze ok. \SI{200}{\celsius}.
  \item \textbf{Sesja M} (\sesjaM) -- napięcie targetu \SI{1500}{V},
        $U=\SIrange{500}{1700}{V}$, prąd w~zakresie
        \SIrange{14}{200}{\micro A}; zdjęcia dla
        $U=\SIrange{500}{1600}{V}$ (arkusz \emph{magnez\_emisja}). Źródło nie było przed tą sesją
        ponownie odgazowane.
\end{itemize}
Dodatkowo wykorzystano serię~2 bez magnesu (10~zdjęć z~odczytami prądu
\SIrange{1}{400}{\micro A}), która posłużyła do sprawdzenia związku jasności
ekranu z~prądem. Wyniki zebrano w~tab.~\ref{tab:emisja-iv}.

\begin{table}[htbp]
  \centering
  \caption{Prąd emisji w~funkcji napięcia w~obu sesjach}
  \label{tab:emisja-iv}
  \begin{tabular}{@{}r r r@{}}
    \toprule
    $U$ [V] & $I_e$ \sesjaB\ [µA] & $I_e$ \sesjaM\ [µA] \\
    \midrule
     500 & 0    & 14 \\
     600 & 0    & 16 \\
     700 & 0    & 14 \\
     800 & 2    & 28 \\
     900 & 3,1  & 45 \\
    1000 & 4    & 69 \\
    1100 & 6    & 118 \\
    1200 & 8    & 180 \\
    1300 & 9    & 190 \\
    1400 & 11,2 & 190 \\
    1500 & 22   & 193 \\
    1600 & --   & 200 \\
    1700 & --   & 200 \\
    \midrule
    $U_\text{target}$ & \SI{2000}{V} & \SI{1500}{V} \\
    \bottomrule
  \end{tabular}
\end{table}

\subsection{Charakterystyki prądowo-napięciowe}

Charakterystyki $I_e(U)$ przedstawiono na rys.~\ref{fig:emisja-iv}.
W~sesji~B prąd pojawia się skokowo przy \SI{800}{V} i~rośnie niemal
liniowo do \SI{1400}{V}; punkt \SI{1500}{V} (\SI{22}{\micro A}, podwojenie
prądu) odbiega od trendu i~może wskazywać na początek wyładowania
lub niestabilności emisji. W~sesji~M dla $U\leq\SI{700}{V}$ płynie
stały prąd ok. \SI{15}{\micro A}, który należy traktować jako prąd
upływu (tło), następnie prąd rośnie wykładniczo do \SI{180}{\micro A}
przy \SI{1200}{V}, po czym nasyca się na poziomie
\SIrange{190}{200}{\micro A}. Przy tym samym napięciu $U$ prąd w~sesji~M
jest ok. 10--20 razy większy niż w~sesji~B (np. \SI{69}{\micro A} wobec
\SI{4}{\micro A} przy \SI{1000}{V}), mimo niższego napięcia targetu.

\subsection{Wpływ warunków źródła na wyniki}

Sesje nie różnią się wyłącznie obecnością magnesu. Przed sesją~B
źródło odgazowano przez wygrzanie w~temperaturze ok.
\SI{200}{\celsius}, natomiast przed sesją~M nie powtórzono tej
procedury. Gazy zaadsorbowane na powierzchni emitera i~elementów
źródła zmieniają pracę wyjścia emitera i~mogą być desorbowane pod
wpływem wiązki, lokalnie zwiększając ciśnienie. Prawdopodobnie
dlatego w~sesji~M obserwuje się:
\begin{itemize}
  \item stały prąd ok. \SI{15}{\micro A} już przy \SI{500}{V}, przed
        właściwym progiem emisji (upływ lub prąd jonowy),
  \item nasycenie prądu powyżej \SI{1200}{V},
  \item rozległą, rozmytą poświatę wokół plamki, która może być
        częściowo świeceniem gazu resztkowego.
\end{itemize}
Wyższy prąd w~sesji~M może wynikać z~dwóch przyczyn, których na
podstawie tych pomiarów nie da się rozdzielić: ogniskowania wiązki
przez magnes (więcej elektronów trafia w~target zamiast w~elektrody
i~ścianki źródła) oraz innego stanu powierzchni emitera i~ciśnienia
w~nieodgazowanym źródle, łącznie z~udziałem prądu upływu. Dlatego
wpływ magnesu oceniono przede wszystkim na podstawie kształtu
i~skupienia plamki.

\begin{figure}[htbp]
  \centering
  \begin{tikzpicture}
    \begin{semilogyaxis}[
      width=0.85\textwidth, height=7cm,
      xlabel={napięcie $U$ [V]}, ylabel={prąd emisji $I_e$ [µA]},
      xmin=450, xmax=1750, ymin=1, ymax=500, xtick distance=200,
      grid=major, grid style={gray!25},
      legend pos=south east, legend cell align=left,
      every axis plot/.append style={thick, mark size=2pt},
      unbounded coords=jump,
    ]
      \addplot[red!60!black, mark=triangle*] table[x=U,y=I_uA,col sep=comma]{\datadir/iv_bez.csv};
      \addlegendentry{\sesjaB\ ($U_\text{t}=\SI{2000}{V}$)}
      \addplot[blue!70!black, mark=*] table[x=U,y=I_uA,col sep=comma]{\datadir/iv_mag.csv};
      \addlegendentry{\sesjaM\ ($U_\text{t}=\SI{1500}{V}$)}
      \draw[gray, dashed] (axis cs:450,15) -- (axis cs:1750,15)
        node[pos=0.2, below, font=\scriptsize] {upływ ok. \SI{15}{\micro A}};
    \end{semilogyaxis}
  \end{tikzpicture}
  \caption{Charakterystyki prądowo-napięciowe emisji w~obu sesjach
    (skala logarytmiczna; punkty z~$I_e=0$ pominięto)}
  \label{fig:emisja-iv}
\end{figure}


\subsection{Metoda analizy zdjęć}

Luminescencja ekranu ma barwę niebieską, a~tło (podświetlenie,
żarnik) -- pomarańczowo-czerwoną. Dla każdego piksela wyznaczono
\emph{nadmiar niebieskiego}
\begin{equation}
  e = \max\bigl(0,\; B - \max(R,G)\bigr),
  \label{eq:nadmiar-niebieskiego}
\end{equation}
gdzie $R,G,B$ to składowe koloru (0--255). Za obszar plamki uznano
piksele z~$e>40$. Na tej podstawie obliczono dwie wielkości:
\begin{itemize}
  \item \textbf{intensywność} $I = \frac{1}{N}\sum e$ (j.u.) -- suma
        nadmiaru niebieskiego znormalizowana do liczby pikseli kadru $N$,
  \item \textbf{pole plamki} -- udział pikseli plamki w~kadrze [\%].
\end{itemize}
Zdjęcia wykonano kamerą z~automatyczną ekspozycją i~przy różnym
kadrowaniu w~każdej sesji, dlatego wartości bezwzględne nie są wprost
porównywalne między sesjami -- porównuje się trendy i~wielkości względne.

\subsection{Intensywność świecenia a~prąd emisji}

Na rys.~\ref{fig:emisja-I} zestawiono intensywność świecenia z~prądem
emisji w~funkcji napięcia. W~sesji~B jasność ekranu rośnie razem
z~prądem (korelacja $r=0{,}99$ dla \SIrange{800}{1100}{V}). W~sesji~M
korelacja wynosi $r=0{,}95$ dla \SIrange{800}{1200}{V}, a~powyżej
\SI{1200}{V} zarówno prąd, jak i~jasność przestają rosnąć -- nasycenie
prądu jest więc rzeczywiste, a~nie wynika tylko z~prześwietlenia
kamery. Warto zauważyć, że w~sesji~M słaba plamka jest widoczna już
przy \SIrange{600}{700}{V}, gdy mierzony prąd nie odbiega jeszcze od
poziomu upływu -- magnes skupia niewielki prąd emisji na małym
obszarze ekranu, przez co staje się on widoczny wcześniej niż bez
magnesu (próg \SI{800}{V}). Dla serii~2 \sesjaB\ korelacja jasności
z~odczytem prądu wynosi $r=0{,}98$ (rys.~\ref{fig:emisja-IvsIe}).

\begin{figure}[htbp]
  \centering
  \begin{tikzpicture}
    \begin{axis}[
      width=0.8\textwidth, height=6.5cm,
      xlabel={napięcie $U$ [V]}, ylabel={intensywność $I$ [j.u.]},
      xmin=450, xmax=1650, ymin=0, ymax=3.5, xtick distance=200,
      grid=major, grid style={gray!25},
      legend pos=north west, legend cell align=left,
      every axis plot/.append style={thick, mark size=2pt},
      axis y line*=left,
    ]
      \addplot[blue!70!black, mark=*] table[x=U,y=I,col sep=comma]{\datadir/mg.csv};
      \addlegendentry{jasność, \sesjaM}
      \addplot[red!60!black, mark=triangle*] table[x=U,y=I,col sep=comma]{\datadir/s1.csv};
      \addlegendentry{jasność, \sesjaB}
      \addlegendimage{blue!70!black, dashed, mark=o}
      \addlegendentry{prąd, \sesjaM\ (oś prawa)}
    \end{axis}
    \begin{axis}[
      width=0.8\textwidth, height=6.5cm,
      xmin=450, xmax=1650, ymin=0, ymax=220,
      axis y line*=right, axis x line=none,
      ylabel={$I_e$ \sesjaM\ [µA]},
    ]
      \addplot[blue!70!black, dashed, mark=o, thick] table[x=U,y=I_uA,col sep=comma]{\datadir/iv_mag.csv};
    \end{axis}
  \end{tikzpicture}
  \caption{Intensywność świecenia ekranu (oś lewa) i~prąd emisji
    w~sesji z~magnesem (oś prawa) w~funkcji napięcia}
  \label{fig:emisja-I}
\end{figure}

\begin{figure}[htbp]
  \centering
  \begin{tikzpicture}
    \begin{loglogaxis}[
      width=0.65\textwidth, height=5.5cm,
      xtick={30,50,100,200,400}, xticklabels={30,50,100,200,400},
      log ticks with fixed point,
      xlabel={prąd emisji $I_e$ [µA] (seria 2, \sesjaB)},
      ylabel={intensywność $I$ [j.u.]},
      grid=major, grid style={gray!25},
      legend pos=north west, legend cell align=left,
    ]
      \addplot[only marks, mark=square*, orange!85!black]
        table[x=etykieta, y=I, col sep=comma,
              restrict expr to domain={\thisrow{I}}{0.01:100}]{\datadir/s2.csv};
      \addlegendentry{pomiar}
      \addplot[domain=30:420, samples=20, black, dashed] {exp(-6.886)*x^1.371};
      \addlegendentry{$I\propto I_e^{1{,}37}$}
    \end{loglogaxis}
  \end{tikzpicture}
  \caption{Intensywność świecenia a~odczyt prądu w~serii~2 \sesjaB}
  \label{fig:emisja-IvsIe}
\end{figure}

\subsection{Skupienie plamki}

Zdjęcia obu sesji przy tych samych napięciach \SIrange{800}{1100}{V}
zestawiono na rys.~\ref{fig:emisja-sbs}, a~parametry plamki w~tab.~\ref{tab:emisja}.
\begin{itemize}
  \item \textbf{Bez magnesu} świeci duży fragment okna ekranu o~ostrych
        krawędziach -- kształt plamki odtwarza geometrię otworu, czyli
        wiązka jest szeroka i~równomiernie oświetla okno.
  \item \textbf{Z~magnesem} świeci zwarty, wydłużony obszar w~części okna,
        otoczony rozległą niebieską poświatą z~promienistymi smugami.
\end{itemize}
Obraz z~sesji~M jest zgodny z~działaniem magnesu jako soczewki
ogniskującej: wiązka zostaje ściśnięta do mniejszego obszaru, zamiast
oświetlać całe okno, dzięki czemu plamka jest widoczna już przy
mniejszym prądzie. Poświata i~smugi mogą pochodzić częściowo
z~rozpraszania światła w~optyce kamery (prześwietlony rdzeń),
a~częściowo z~elektronów poza głównym ogniskiem lub ze świecenia
gazu resztkowego w~nieodgazowanym źródle. Ze względu na brak
skali na zdjęciach i~inne kadrowanie nie można podać średnicy plamki
w~mikrometrach.

\begin{table}[htbp]
  \centering
  \caption{Parametry emisji i~plamki w~obu sesjach}
  \label{tab:emisja}
  \begin{tabular}{@{}lccc@{}}
    \toprule
    Parametr & M (\sesjaM) & B (\sesjaB) & seria 2 (\sesjaB) \\
    \midrule
    Odgazowanie źródła & nie & tak (\SI{200}{\celsius}) & -- \\
    Napięcie targetu & \SI{1500}{V} & \SI{2000}{V} & -- \\
    Próg emisji (prąd) & \SI{800}{V} & \SI{800}{V} & -- \\
    Pierwsza widoczna plamka & \SI{600}{V} & \SI{800}{V} & -- \\
    Prąd maksymalny & \SI{200}{\micro A} & \SI{22}{\micro A} & \SI{400}{\micro A} \\
    Maks. intensywność $I$ [j.u.] & 3,19 & 1,67 & 2,89 \\
    Maks. pole plamki [\% kadru] & 3,98 & 2,07 & 3,91 \\
    \bottomrule
  \end{tabular}
\end{table}

% ---------------- zestawienie side-by-side ----------------
\newcommand{\sbs}[2]{%
  \begin{subfigure}[t]{0.28\textwidth}\centering
    \includegraphics[width=\linewidth]{\figdir/#1}\caption{#2}
  \end{subfigure}}

\begin{figure}[p]
  \centering
  \makebox[0.28\textwidth]{\textbf{\sesjaM}}\hspace{1.5em}\makebox[0.28\textwidth]{\textbf{\sesjaB}}\par\medskip
  \sbs{mg_03.jpg}{\SI{800}{V}, \SI{28}{\micro A}}\hspace{1.5em}
  \sbs{s1_02.jpg}{\SI{800}{V}, \SI{2}{\micro A}}\par\medskip
  \sbs{mg_04.jpg}{\SI{900}{V}, \SI{45}{\micro A}}\hspace{1.5em}
  \sbs{s1_03.jpg}{\SI{900}{V}, \SI{3,1}{\micro A}}\par\medskip
  \sbs{mg_05.jpg}{\SI{1000}{V}, \SI{69}{\micro A}}\hspace{1.5em}
  \sbs{s1_04.jpg}{\SI{1000}{V}, \SI{4}{\micro A}}\par\medskip
  \sbs{mg_06.jpg}{\SI{1100}{V}, \SI{118}{\micro A}}\hspace{1.5em}
  \sbs{s1_05.jpg}{\SI{1100}{V}, \SI{6}{\micro A}}
  \caption{Plamka luminescencji \sesjaM\ (lewa kolumna) i~\sesjaB\ (prawa
    kolumna) przy tych samych napięciach -- zbliżenia okna ekranu.
    Z~magnesem świeci zwarty obszar z~poświatą, bez magnesu -- szeroki
    fragment okna o~ostrych krawędziach}
  \label{fig:emisja-sbs}
\end{figure}

\begin{figure}[htbp]
  \centering
  \begin{subfigure}[t]{0.28\textwidth}\centering
    \includegraphics[width=\linewidth]{\figdir/mg_02.jpg}\caption{\SI{700}{V}}\end{subfigure}\hspace{1.5em}
  \begin{subfigure}[t]{0.28\textwidth}\centering
    \includegraphics[width=\linewidth]{\figdir/mg_09.jpg}\caption{\SI{1400}{V}}\end{subfigure}\hspace{1.5em}
  \begin{subfigure}[t]{0.28\textwidth}\centering
    \includegraphics[width=\linewidth]{\figdir/mg_11.jpg}\caption{\SI{1600}{V}}\end{subfigure}
  \caption{Sesja \sesjaM: plamka widoczna przy prądzie na poziomie upływu
    (\SI{700}{V}) oraz w~zakresie nasycenia prądu (\SIrange{1400}{1600}{V})}
  \label{fig:emisja-mg-skrajne}
\end{figure}

\begin{figure}[htbp]
  \centering
  \begin{subfigure}[t]{0.32\textwidth}\centering
    \includegraphics[width=\linewidth]{\figdir/mg_pelny.jpg}\caption{\sesjaM, \SI{1400}{V}}\end{subfigure}\hfill
  \begin{subfigure}[t]{0.32\textwidth}\centering
    \includegraphics[width=\linewidth]{\figdir/s1_pelny.jpg}\caption{\sesjaB, \SI{1100}{V}}\end{subfigure}\hfill
  \begin{subfigure}[t]{0.32\textwidth}\centering
    \includegraphics[width=\linewidth]{\figdir/s2_pelny.jpg}\caption{\sesjaB, seria 2}\end{subfigure}
  \caption{Pełne kadry przy maksymalnej emisji w~każdej serii}
  \label{fig:emisja-pelne}
\end{figure}

\subsection{Wnioski}

\begin{enumerate}
  \item W~obu sesjach próg emisji wynosi ok. \SI{800}{V}.
  \item Z~magnesem plamka jest zwarta i~widoczna już przy
        \SI{600}{V}--\SI{700}{V}, mimo że prąd nie odbiega jeszcze od
        poziomu upływu; bez magnesu świeci szeroko całe okno dopiero od
        \SI{800}{V}. Wskazuje to na ogniskujące działanie pola
        magnetycznego.
  \item W~sesji z~magnesem prąd nasyca się przy ok.
        \SI{200}{\micro A} powyżej \SI{1200}{V}, co potwierdza również
        przestający rosnąć obraz plamki.
  \item Jasność ekranu silnie koreluje z~prądem emisji
        ($r=0{,}95$--$0{,}99$) i~może służyć jako jego względna miara.
  \item Bezpośrednie porównanie wartości prądów między sesjami jest
        ograniczone: źródło było odgazowane tylko przed sesją bez magnesu,
        sesje różnią się napięciem targetu (\SI{2000}{V} vs
        \SI{1500}{V}). Z~magnesem prąd jest 10--20 razy większy, co może
        wynikać zarówno z~ogniskowania, jak i~ze stanu źródła. Prąd
        upływu, nasycenie i~poświata w~sesji z~magnesem wynikają
        prawdopodobnie z~braku odgazowania.
  \item Do ilościowej oceny średnicy plamki potrzebne są: stały kadr
        i~ekspozycja kamery oraz wzorzec skali w~polu widzenia.
\end{enumerate}
